\chapter{阿代尔群的自守表示II：若干性质与局部理论}\label{L5}
\lecturedate{2020/03/02}
\sourcepage{36}
本讲主要沿用\bref{L5refBG}{BG}的第7.2、7.3节。
\section{\texorpdfstring{$(\g,K)\times G(\A^f)$-模}{(g,K)×G(Aᶠ)-模}}\label{L5S1}
我们已经定义了空间$\mathcal A(G,\omega)$和$\mathcal A_0(G,\omega)$，$(\g,K)\times G(\A^f)$作用于这两个空间。$K$的无穷分量是$G(\R)$的极大紧子群$O$。$\g_0$为$G(\R)$的实Lie代数，$\g$为其复化。记$\mathfrak{o}_0\subseteq\g_0$为$O$的实Lie代数。注意，$G(\R)$在$\g$上有伴随作用$\operatorname{Ad}(\cdot)$。
% [Tutor Note: 原书 p36 的小节标题与第二句均作 $(\g,K)\times G(\A)$；实际作用于 $\mathcal A(G,\omega)$ 的是 $(\g,K)\times G(\A^f)$（$K$ 通过其无穷分量 $O$ 起作用，有限分量已含在 $G(\A^f)$ 中），译文按 $G(\A^f)$ 处理。]
\originalnumber{1}
\begin{definition}{$(\g,O)$-模及其整体版本}\label{L5S1O1}
\begin{enumerate}
\item 一个\textbf{$(\g,O)$-模}\znote{$(\g,O)$-模又称 Harish-Chandra 模：它把 $G(\R)$ 的表示拆成两片，一片是 Lie 代数 $\g$ 的代数作用（无穷小方向），另一片是极大紧子群 $O$ 的群作用。$O$-有限性保证任取一个向量，它张成的 $O$-子模有限维，这正是「可容许」能逐项计数的前提。当 $K=O\times K_f$ 时，$(\g,O)$-模与 $G(\A^f)$ 的光滑作用合起来即 $(\g,K)\times G(\A^f)$-模。[来源：Bernstein--Gelbart, \textit{An introduction to the Langlands Program}, 第 7 章（本讲参考文献）。]}是带有$\g$及$O$的作用的复向量空间$V$（也可记为$(\pi,V)$，用$\pi$表示这些作用），满足：
\begin{enumerate}[label=(\arabic*)]
\item $V$中的所有向量都是$O$-有限的；
\item 对$X\in\mathfrak{o}_0$、$v\in V$，$X\cdot v=\left.\frac d{dt}(\exp(tX)\cdot v)\right|_{t=0}$；
% [Tutor Note: 原书 p36 定义 1.1(2) 作 $X\in\mathfrak{a}_0$，但上文已记 $\mathfrak{o}_0$ 为 $O$ 的实 Lie 代数，$\mathfrak{a}_0$ 无定义；译文按 $\mathfrak{o}_0$ 处理。]
\item $\g$与$O$的作用相容，即对所有$y\in O$、$X\in\g$，有$y\cdot(X\cdot v)=(\operatorname{Ad}(y)X)\cdot(y\cdot v)$。
\end{enumerate}
\item 一个\textbf{$(\g,O)\times G(\A^f)$-模}是带有$G(\A^f)$的光滑作用的$(\g,O)$-模，且该作用与$(\g,O)$的作用交换。这里光滑是指：每个向量$v\in V$都被$G(\A^f)$的某个紧开子群$U$固定。
\end{enumerate}
\end{definition}
\sourcepage{37}
$\mathcal A(G):=\mathcal A(G,\omega)$与$\mathcal A_0(G):=\mathcal A_0(G,\omega)$都是$(\g,O)\times G(\A^f)$-模。实Lie代数$\g_0$在$\mathcal A(G)$上的作用为
\[X\cdot f(\alpha_\infty)=\left.\frac d{dt}\{f(\alpha_\infty\exp(tX))\}\right|_{t=0}.\]
复化Lie代数$\g$的作用由$(X+iY)\cdot f=X\cdot f+i(Y\cdot f)$定义，其中$X,Y\in\g_0$。
\originalnumber{2}
\begin{definition}{可容许性与不可约性}\label{L5S1O2}\idx{可容许表示}\idx{不可约表示}
\begin{enumerate}
\item 若$K$的每个不可约表示$\sigma$在$V$中的重数都有限，就称$(\g,O)\times G(\A^f)$-模$(\pi,V)$\textbf{可容许}\znote{可容许性要求每个不可约表示 $\sigma$ 在 $V$ 中只出现有限次，于是 $V$ 可按 $K$-类型分解为代数直和。缺了这一条，$V$ 可能含无限重数，迹公式与 $L$ 函数的局部因子都无从定义。局部理论里「$V^U$ 有限维」是同一条件的另一种说法；整体情形则由无穷处的 $O$-有限性与有限处的光滑性拼合而成。[来源：Bernstein--Gelbart, \textit{An introduction to the Langlands Program}, 第 7 章（本讲参考文献）。]}。
\item 若$(\pi,V)$没有在$\g$、$O$和$G(\A^f)$的作用下保持不变的非零真子空间，就称它\textbf{不可约}。
\end{enumerate}
\end{definition}
有如下对应的局部定义。
\originalnumber{3}
\begin{definition}{局部可容许表示}\label{L5S1O3}
\begin{description}
\item[(A)] 当$l$为阿基米德素位时，若$O$的每个不可约表示$\sigma$在$V_l$中的重数都有限，就称$(\g,O)$-模$(\pi_l,V_l)$可容许。若它没有在$\g$和$O$作用下不变的非零真子空间，则称为不可约。
\item[(NA)] 当$l$为非阿基米德素位时，若$G(F_l)$的光滑表示$(\pi_l,V_l)$满足：对任意开子群$U\subseteq G(F_l)$，$V_l^U$有限维，就称它可容许。若它没有非零真不变子空间，则称为不可约。此外，若不可约可容许表示$(\pi_l,V_l)$的$G(\OO_l)$-不变向量空间$V_l^{G(\OO_l)}$非零，就称它\textbf{非分歧}，或称\textbf{球面表示}；这里$\OO_l$为$F_l$的赋值环。
\end{description}
\end{definition}
\idx{非分歧表示}\idx{球面表示}
\originalnumber{4}
\begin{note}[L5S1O4]
已知，对于有限素数$p$，若$(\pi_p,V_p)$是$G(\Q_p)$的不可约可容许表示，则对于这里使用的标准极大紧子群，有$\dim V_p^{G(\Z_p)}\leq1$。因此，它非分歧$\iff\dim V_p^{G(\Z_p)}=1$。
\end{note}
现在可以精确定义不可约尖点自守表示。
\originalnumber{5}
\begin{definition}{尖点自守表示}\label{L5S1O5}
$G(\A)$的（不可约）\textbf{尖点自守表示}，是出现在$\mathcal A_0(G)$中的一个（不可约）可容许$(\g,O)\times G(\A^f)$-模。
\end{definition}
\sourcepage{38}
\section{自守表示的若干性质}\label{L5S2}
下面的因子分解定理由Flath给出，使我们能够通过$G(\Q_l)$的局部表示来研究$G(\A)$的自守表示。
\originalnumber{1}
\begin{theorem}\label{L5S2O1}\idx{限制张量积}
$G(\A)$的不可约可容许表示$(\pi,V)$是一个限制张量积$\bigotimes_l'V_l$\znote{限制张量积是无穷张量积的收敛版本：只取那些在几乎所有素位处等于指定基向量 $u_p^0$ 的纯张量，再作有限线性组合。这样得到的空间不必是整个张量积，却仍承载各局部表示的作用。Flath 定理因此把整体表示彻底拆成局部数据，使整体问题可以逐素位处理，这正是 Langlands 纲领「局部到整体」的技术支点。[来源：Bernstein--Gelbart, \textit{An introduction to the Langlands Program}, 第 7 章（本讲参考文献）。]}，相对于一族向量$(u_p^0)$定义；在几乎所有$p$处，$u_p^0$是$V_p^{G(\Z_p)}$的基向量。各$V_l$不可约可容许，且几乎所有$V_p$非分歧。换言之，$V$的每个元素都是形如$u=\bigotimes_lu_l$的纯张量的有限线性组合，其中$u_l\in V_l$，且对几乎所有$p$都有$u_p=u_p^0$。
\end{theorem}
\originalnumber{2}
\begin{corollary}\label{L5S2O2}
不可约自守表示$\pi$有一个限制张量积分解$\pi\cong\bigotimes_l'\pi_l$，其中对几乎所有$p$，$\pi_p^{G(\Z_p)}\ne0$，即$\pi_p$非分歧。
\end{corollary}
现在集中讨论$\mathcal A_0(G)$。事实上，由Langlands的一个定理，尖点表示是基本对象：其他自守表示都可通过抛物诱导由它们构造出来。

取$G=\GL_n$，于是$O=\operatorname O(n)$。关于$\mathcal A_0(G)$的主要结论如下。
\originalnumber{3}
\begin{theorem}\label{L5S2O3}\idx{重数一定理}\idx{强重数一定理}
\begin{enumerate}[label=(\roman*)]
\item 空间$\mathcal A_0(G)$是$G(\A)$的不可约自守表示的代数直和：
\[\mathcal A_0(G)=\bigoplus_\pi\mathcal A_0(\pi).\]
这里$\mathcal A_0(\pi)$是与$\pi$同构的直和项。
\item （重数一）$G(\A)$的一个不可约表示在$\mathcal A_0(G)$中的重数至多为$1$。换言之，若$\pi_1,\pi_2$是其中的两个不可约表示，则
\[\pi_1\cong\pi_2\iff\mathcal A_0(\pi_1)=\mathcal A_0(\pi_2).\]
\item （强重数一）若$\mathcal A_0(\pi_1),\mathcal A_0(\pi_2)\subseteq\mathcal A_0(G)$，且几乎所有$l$处的$\pi_{1,l}$与$\pi_{2,l}$同构，则它们作为$\mathcal A_0(G)$的子空间相同：$\mathcal A_0(\pi_1)=\mathcal A_0(\pi_2)$。
\end{enumerate}
\end{theorem}
\originalnumber{4}
\begin{note}[L5S2O4]
当$G=\SL_n$、$n\geq3$时，“重数一”性质不成立。
\end{note}
\sourcepage{39}
\section{局部理论}\label{L5S3}
本节设$F$为非阿基米德局部域，$\ord(\cdot)$为其规范化离散赋值。设$\OO$为$F$中的离散赋值环，$\Pp$为$\OO$的唯一非零素理想，剩余域$\OO/\Pp$有$q$个元素。于是$F$上的绝对值为$|x|=q^{-\ord(x)}$。

令$G=\GL_2(F)$，$K=\GL_2(\OO)$为$G$的一个极大紧子群。令$B=NM$，其中
\[M=\left\{m=m(a)=\begin{pmatrix}a_1&0\\0&a_2\end{pmatrix}\ \middle|\ a_1,a_2\in F^\times\right\},\qquad N=\left\{n=n(x)=\begin{pmatrix}1&x\\0&1\end{pmatrix}\ \middle|\ x\in F\right\}.\]
由Iwasawa分解，$G=BK$。记$Z$为$G$的中心，则$Z\cong F^\times$。

对$F^\times$的一对光滑特征标$\mu=(\mu_1,\mu_2)$，定义表示$I(\mu)$如下：
\begin{equation}\tag{3.1}\label{L5eq31}
I(\mu)=\left\{f:G\longrightarrow\C\ \middle|\ \begin{array}{l}f\text{ 光滑，即在某个开子群 }U_f\subseteq G\\\text{的右平移下不变；}\\f(bg)=\mu(b)\delta(b)^{1/2}f(g),\quad b\in B,\ g\in G\end{array}\right\}.
\end{equation}
这里，对$b=n\,m(a)$，$\mu(b)=\mu_1(a_1)\mu_2(a_2)$，$\delta(b)=|a_1/a_2|=q^{-\ord(a_1)+\ord(a_2)}$。
\originalnumber{1}
\begin{note}[L5S3O1]
若$f$被开子群$U_f$固定，则对任意$g\in G$和$u\in U_f$，
\[(g^{-1}ug)\cdot(g^{-1}\cdot f)(x)=f(xg^{-1}u)=f(xg^{-1}).\]
因此，平移后的函数$g^{-1}\cdot f$被$g^{-1}U_fg$固定。对$f$本身，可以把固定它的开子群替换为$U_f\cap K$，从而假定$U_f\subseteq K$。
\end{note}
$I(\mu)$的中心特征标是$\omega=\mu_1\mu_2$，即对$z\in Z$、$g\in G$，有$f(zg)=\omega(z)f(g)$。定义特征标$\omega_s(x)=|x|^s=q^{-s\ord(x)}$（$x\in F^\times$），则$\delta(b)=\omega_1(a_1)\omega_{-1}(a_2)$。有
\[\begin{cases}\mu(b)\delta(b)^{1/2}=\mu_1(a_1)\mu_2(a_2)\omega_{1/2}(a_1/a_2),&\text{若 }\mu_1\mu_2^{-1}=\omega_1,\\\mu(b)\delta(b)^{1/2}=\mu_1(a_1)\mu_2(a_2)\omega_{-1/2}(a_2/a_1),&\text{若 }\mu_1\mu_2^{-1}=\omega_{-1}.\end{cases}\]
% [Tutor Note: 原书 p39 第二行印作 $\omega_{-1/2}(a_1/a_2)$，与第一行相同；由 $\delta(b)=|a_1/a_2|$ 及条件 $\mu_1\mu_2^{-1}=\omega_{-1}$，该项应为 $\omega_{-1/2}(a_2/a_1)$，译文已按正确写法处理。]
\sourcepage{40}
\originalnumber{2}
\begin{theorem}\label{L5S3O2}\idx{主系列表示}\idx{特殊表示}
\begin{enumerate}[label=(\arabic*)]
\item $I(\mu)$不可约，当且仅当$\mu_1\mu_2^{-1}\ne\omega_{\pm1}$。
\item 若$\mu_1\mu_2^{-1}=\omega_1$，则$I(\mu)$有一个一维商，$G$通过特征标$\chi\circ\det$作用于它，其中$\chi=\mu_1\omega_{-1/2}=\mu_2\omega_{1/2}$；另有一个无限维不可约子表示$\sigma(\mu)$。
\item 若$\mu_1\mu_2^{-1}=\omega_{-1}$，则$I(\mu)$有一个一维子模，$G$通过特征标$\chi\circ\det$作用于它，其中$\chi=\mu_1\omega_{1/2}=\mu_2\omega_{-1/2}$；另有一个无限维不可约商$\sigma(\mu)$。
\item 这些表示之间的等价关系只有如下两种。令$\mu'=(\mu_2,\mu_1)$：
\begin{itemize}
\item 若$\mu_1\mu_2^{-1}\ne\omega_{\pm1}$，则$I(\mu)\cong I(\mu')$；
\item 若$\mu_1\mu_2^{-1}=\omega_{\pm1}$，则$\sigma(\mu)\cong\sigma(\mu')$。
\end{itemize}
\end{enumerate}
\end{theorem}
\originalnumber{3}
\begin{note}[L5S3O3]
若$\mu_1\mu_2^{-1}=\omega_{-1}$，则$G$通过特征标$\chi\circ\det$作用于$I(\mu)$的一维子模$\C(\chi\circ\det)$。
\end{note}
若$G$的不可约可容许表示$(\pi,V)$满足：对所有$\mu$，$\Hom_G(\pi,I(\mu))=0$，就称其为\textbf{超尖点表示}\znote{超尖点表示是 $\GL_2(F)$ 最「原子」的表示：它与任何主系列都不相交，既不出现在诱导表示里，也不能由更小的表示诱导得到；其余不可约表示都可由主系列及其子商刻画，所以它是构造全部表示的基本砖块。等价的判别是矩阵系数在 $Z(F)N(F)\backslash G(F)$ 上紧支，而定义里的 $\Hom_G(\pi,I(\mu))=0$ 最便于验证。[来源：Bernstein--Gelbart, \textit{An introduction to the Langlands Program}, 第 7 章（本讲参考文献）。]}。因此，超尖点表示、不可约主系列表示\znote{主系列表示指从 Borel 子群 $B=NM$ 诱导的表示：式 (3.1) 用 $f(bg)=\mu(b)\delta(b)^{1/2}f(g)$ 定义 $I(\mu)$，其中 $\delta$ 为模函数，$1/2$ 次幂使对偶与酉性对称。$\GL_2$ 的不可约可容许表示只有四类：超尖点、主系列、特殊表示、一维表示。主系列一般可约；可约时它的子商正是特殊表示与一维表示。[来源：Bernstein--Gelbart, \textit{An introduction to the Langlands Program}, 第 7 章（本讲参考文献）。]}$I(\mu)$、特殊表示$\sigma(\mu)$和一维表示$\chi\circ\det$，在同构意义下给出了$G$的所有不可约可容许表示。\idx{超尖点表示}

对于局部域$F$，若特征标$\omega:F^\times\longrightarrow\C^\times$满足$\omega(\OO^\times)=\{1\}$，则称它\textbf{非分歧}。非分歧特征标形如$\omega(x)=t^{\ord(x)}$，其中$t\in\C^\times$。若$\mu_1,\mu_2$都非分歧，则$\mu=(\mu_1,\mu_2)$在$B\cap K$上的限制平凡。限制到$K$给出同构
% [Tutor Note: 原书 p40 作 $\mu=(\mu_1,\mu_1)$ 在 $B\cap K$ 上平凡；此处应为 $\mu=(\mu_1,\mu_2)$，译文已按正确写法处理。]
\[I(\mu)\cong C^\infty((B\cap K)\backslash K),\]
其中常值函数的原像恰好是$I(\mu)^K$。因此$\dim I(\mu)^K=1$，它有非零球面向量。对每一对$(t_1,t_2)\in(\C^\times)^2$，都有一个表示$I(\mu)$，其中$\mu=(t_1^{\ord(\cdot)},t_2^{\ord(\cdot)})$；由\eref{定理}{L5S3O2}，它不可约当且仅当$t_1t_2^{-1}\ne q^{\pm1}$。容易看出，对每个$(t_1,t_2)\in(\C^\times)^2$，存在不可约可容许非分歧表示$\pi(t_1,t_2)$，使
% [Tutor Note: 原书 p40 末作「对每个 $(t_1,t_2)\in(\A^\times)^2$」，但上下文是 $\C^\times$ 中的参数，$\A^\times$ 在此无意义；译文按 $(\C^\times)^2$ 处理。]
\begin{enumerate}[label=(\arabic*)]
\item 若$t_1t_2^{-1}\ne q^{\pm1}$，则$\pi(t_1,t_2)\cong\pi(t_2,t_1)=I(\mu)$，其中$\mu=(t_1^{\ord(\cdot)},t_2^{\ord(\cdot)})$。
\sourcepage{41}
\item 若$t_1t_2^{-1}=q^{\pm1}$，则$\pi(t_1,t_2)=\C(\chi\circ\det)$，为$I(\mu)$的一维成分。
\end{enumerate}
\originalnumber{4}
\begin{theorem}\label{L5S3O4}
$G$的每个不可约可容许非分歧表示都同构于某个$\pi(t_1,t_2)$。
\end{theorem}
\originalnumber{5}
\begin{corollary}\label{L5S3O5}
$G$的不可约可容许非分歧表示的同构类，与复群$G^\vee:=\GL_2(\C)$中对角矩阵$\begin{pmatrix}t_1&0\\0&t_2\end{pmatrix}$的共轭类之间存在双射，具体为
\[\pi=\pi(t_1,t_2)\longleftrightarrow t(\pi):=\begin{pmatrix}t_1&0\\0&t_2\end{pmatrix}\text{ 的共轭类}.\]
\end{corollary}
\eref{推论}{L5S3O5}中的半单共轭类$t(\pi)$称为$\pi$的\textbf{Satake参数}\znote{Satake 参数把非分歧表示翻译成对偶群中的半单共轭类：$\pi(t_1,t_2)$ 对应 $\GL_2(\C)$ 中 $\operatorname{diag}(t_1,t_2)$ 的共轭类。称「无序对」，是因为交换 $t_1,t_2$ 给出同构的表示。这是 Satake 同构在 $\GL_2$ 的特例：非分歧表示的分类等价于对偶环面在 Weyl 群作用下的轨道分类。[来源：Bernstein--Gelbart, \textit{An introduction to the Langlands Program}, 第 7 章（本讲参考文献）。]}；有时也把无序对$(t_1,t_2)$称为Satake参数。\idx{Satake参数}

对整数$r\geq0$，定义
\[K_0(r)=\left\{\begin{pmatrix}a&b\\c&d\end{pmatrix}\in\GL_2(\OO)\ \middle|\ \ord(c)\geq r\right\}.\]
设$\omega$为$F^\times$的光滑特征标。其\textbf{导子指数}$c(\omega)$为使$\omega$在$U^N$上平凡的最小非负整数$N$，其中$U^0=\OO^\times$，$U^n=1+\Pp^n$（$n\geq1$）；于是$\omega$在所有$n\geq N$的$U^n$上都平凡。对$r\geq\max\{1,c(\omega)\}$，在$K_0(r)$上定义同样记为$\omega$的特征标$\omega(k)=\omega(a)$，其中$k=\begin{pmatrix}a&b\\c&d\end{pmatrix}$。若$c(\omega)=0$，则$r=0$时取$K_0(0)$上的平凡特征标。下面的结论由Casselman给出。\idx{导子指数}
% [Tutor Note: 原书 p41 的 $K_0(r)$ 定义未写明矩阵属于 $\GL_2(\OO)$；导子指数的定义原书作「$\omega$ 在 $1+\Pp^n$ 上平凡（$\forall n\geq N$）」，$N=0$ 时 $1+\Pp^0=\OO$，应为 $U^0=\OO^\times$；在 $K_0(r)$ 上定义 $\omega(k)=\omega(a)$ 时原书亦未写出 $r\geq\max\{1,c(\omega)\}$ 与 $c(\omega)=0$ 时取平凡特征标这两个条件。译文均按正确写法处理。]
\originalnumber{6}
\begin{theorem}\label{L5S3O6}
设$(\pi,V)$为$G$的无限维不可约可容许表示，中心特征标为$\omega_\pi$。则存在最小整数$c(\pi)\geq c(\omega_\pi)$，使对所有$r\geq c(\pi)$，有
\[V^{K_0(r),\omega_\pi}:=\{v\in V\mid\pi(k)v=\omega_\pi(k)v,\ k\in K_0(r)\}\ne0.\]
此外，$\dim V^{K_0(c(\pi)),\omega_\pi}=1$。
\end{theorem}
\eref{定理}{L5S3O6}中的整数$c(\pi)$称为$\pi$的\textbf{导子指数}。
\originalnumber{7}
\begin{note}[L5S3O7]
若$\pi$非分歧，则$c(\pi)=0$。若$\pi\cong\sigma(\mu)$，其中$\sigma(\mu)$如\eref{定理}{L5S3O2}所定义，且$\mu_1,\mu_2$非分歧，则$c(\pi)=1$；这些表示称为非分歧特殊表示。在中心特征标非分歧的情形下，其余无限维不可约可容许表示均满足$c(\pi)\geq2$；特别地，所有超尖点表示的导子指数至少为$2$。
% [Tutor Note: 原书 p41 的 Remark 3.7 称「$c(\pi)\geq2$ 对其余所有无限维不可约可容许表示」，该断言只在中心特征标非分歧时成立（否则 $c(\pi)$ 还受 $c(\omega_\pi)$ 影响）；译文已补此前提出。]
\end{note}
\sourcepage{42}
\section{阿基米德理论}\label{L5S4}
本节取$G=\GL_2(\R)$，$\g_0=\operatorname{Lie}(G)$，$\g=\g_0\otimes_{\R}\C$。固定极大紧子群为$\operatorname O(2)$。对整数$n\geq2$，有不可约$(\g,\operatorname O(2))$-模$\operatorname{DS}_n=(\pi_n,V)$，定义如下：
\[V=\bigoplus_{\substack{|l|\geq n\\l\equiv n\pmod2}}\C v_l,\]
其中$\operatorname O(2)$的作用为
\[\pi_n\!\left(\begin{pmatrix}\cos\theta&\sin\theta\\-\sin\theta&\cos\theta\end{pmatrix}\right)v_l=e^{il\theta}v_l,\qquad\pi_n\!\left(\begin{pmatrix}1&0\\0&-1\end{pmatrix}\right)v_l=v_{-l}.\]
令
\[X_\pm=\frac12\begin{pmatrix}1&\pm i\\\pm i&-1\end{pmatrix}\in\g,\qquad H=\begin{pmatrix}0&-i\\i&0\end{pmatrix}\in\g.\]
$X_\pm,H$构成$\mathfrak{sl}_2(\C)\subseteq\g$的一组基，其Killing型非退化。\textbf{Casimir元素}\idx{Casimir元素}
\[C:=H^2+2X_+X_-+2X_-X_+\]
生成$U(\mathfrak{sl}_2(\C))$的中心。令
% [Tutor Note: 原书 p42 作「Casimir 元素生成中心 $\mathfrak{Z}(\mathfrak{g})$」；$C\in U(\mathfrak{g})$ 而非 $\mathfrak{g}$，该记号有歧义（若指 Lie 代数中心则不确），译文按 $U(\mathfrak{sl}_2(\C))$ 的中心处理。]
\[\pi_n(X_\pm)v_l=\frac12(n\pm l)v_{l\pm2},\]
于是$\pi_n(X_-)v_n=0$、$\pi_n(X_+)v_{-n}=0$。最后令$\pi_n(C)v=n(n-2)v$。
\begin{exercise*}\label{L5S4EX}
验证：$(\pi_n,V)$给出了一个不可约$(\g,\operatorname O(2))$-模。
\end{exercise*}
\lref{L6}将利用局部理论研究$(\g,O)\times G(\A^f)$-模，并建立$\mathcal A_0(G)$与尖点模形式空间之间的联系。
\needspace{5\baselineskip}\section{第5讲参考文献}\label{L5S5}
\begin{references}\refitem{L5refBG}{BG} J. Bernstein and S. Gelbart, eds, \textit{An introduction to the Langlands Program}, Birkh\"auser, 2004, Chapter 7 written by Stephen S. Kudla.\end{references}
